\documentclass[11pt,a4paper]{article}
\usepackage[T1]{fontenc}
\usepackage[utf8]{inputenc}
\usepackage[margin=1in]{geometry}
\usepackage{booktabs}
\usepackage{longtable}
\usepackage{array}
\usepackage{amsmath}
\usepackage{hyperref}
\usepackage{graphicx}

\title{GlycoSMILES2BAP: Automated Stereochemistry-Preserving Glycan Input Preparation for Standalone AlphaFold 3}
\author{Qiang Xia\\
\textit{Zhejiang Xinghe Tea Technology Co., Ltd., Hangzhou, Zhejiang, China}\\
\texttt{xiaqiang@xinghetea.com}}
\date{}

\begin{document}
\maketitle

\begin{abstract}
\noindent\textbf{Motivation:} AlphaFold 3 achieves unprecedented accuracy for protein-glycan complex structure prediction, but Huang et al. (2025) identified a critical limitation: standard input formats produce stereochemically incorrect glycan structures. The only format preserving stereochemistry, CCD+bondedAtomPairs (BAP), requires manual atom-by-atom specification taking 30--60 minutes per structure, creating a prohibitive barrier for large-scale glycan modeling.

\noindent\textbf{Results:} We present GlycoSMILES2BAP, an automated pipeline converting standard glycan notations (IUPAC-condensed, WURCS) to standalone AF3-compatible CCD+BAP format. The pipeline employs three mechanism-driven modules: (1) a CCD mapper with anomeric position tracking that correctly handles sialic acids (C2 anomeric carbon) and pentoses (O4 ring oxygen), (2) a topology parser extracting linkage information from branched structures, and (3) a BAP generator producing explicit atom-pair bond specifications. On a benchmark of 50 diverse glycan structures, GlycoSMILES2BAP achieves 98.5\% epimer accuracy, 98.2\% anomeric accuracy, and 96.8\% linkage accuracy, with bootstrap 95\% confidence intervals of [96.2\%, 99.8\%], [95.8\%, 99.6\%], and [93.5\%, 99.2\%], respectively, compared to $\sim$60\% for SMILES-based approaches. Processing time is $0.82 \pm 0.15$ seconds per structure versus 30--60 minutes manually. Ablation studies confirm each module's contribution: removing CCD mapping reduces epimer accuracy by 15.5 percentage points; removing anomeric tracking causes sialic acid failures. Validation against 10 literature-reported test cases shows 100\% correction within that set, and representative database-scale processing achieves 94\% success on 100 GlyTouCan structures. The present study addresses AF3 input construction rather than all downstream prediction limitations for complex polysaccharides.

\noindent\textbf{Conclusions:} GlycoSMILES2BAP bridges the gap between accessible glycan notations and AF3's stereochemistry-preserving input format, enabling accurate, reproducible AF3 input preparation without manual specification overhead. Public AlphaFold Server support for glycan-containing inputs remains limited, so end-to-end validation for complex polysaccharides still depends on standalone AF3 execution and glycan-aware structural assessment.

\noindent\textbf{Availability:} \url{https://github.com/ShawnXiha/glycobap}

\noindent\textbf{Keywords:} AlphaFold 3, glycans, stereochemistry, structure prediction, CCD, bondedAtomPairs
\end{abstract}

\section{Introduction}

Glycosylation is one of the most prevalent post-translational modifications in eukaryotic proteins, affecting over 50\% of the human proteome (Varki, 2017). Glycans play essential roles in protein folding, cell-cell recognition, immune responses, and pathogen-host interactions (Helenius and Aebi, 2004). The structural diversity of glycans, arising from variations in monosaccharide composition, linkage positions, anomeric configurations, and branching patterns, poses significant challenges for structural characterization. GlyTouCan, the international glycan structure repository, catalogs over 200,000 unique entries (Tiemeyer et al., 2024), reflecting the substantial scale of glycan diversity that researchers must navigate.

AlphaFold 3 (AF3) has transformed structural biology by achieving unprecedented accuracy in predicting protein-ligand complex structures, including glycosylated proteins (Abramson et al., 2024). This breakthrough has generated considerable interest in the glycobiology community, as accurate glycan structure prediction could accelerate research in glycoprotein engineering, vaccine design, and therapeutic development. However, a fundamental barrier prevents researchers from fully leveraging AF3's capabilities for glycan modeling.

Recent systematic analysis by Huang et al. (2025) revealed that AF3's standard input formats systematically fail to preserve glycan stereochemistry. Their evaluation identified a stark accuracy gradient: the widely used SMILES format achieves only $\sim$62\% stereochemistry accuracy due to epimer confusion and anomeric inversion. The userCCD format improves to $\sim$82\% accuracy but still introduces linkage position errors in complex branched structures.

Critically, only the CCD+bondedAtomPairs (BAP) format achieves near-perfect accuracy ($\sim$100\%) by explicitly specifying atom-to-atom bonds. However, this format requires researchers to manually identify and annotate each glycosidic bond, a process that takes 30--60 minutes per structure for an expert. For a modest library of 100 glycan variants, manual BAP specification would require 50--100 hours of expert time, effectively prohibiting large-scale glycan structure prediction.

We present GlycoSMILES2BAP, an automated pipeline that bridges this gap by converting standard glycan notations (IUPAC-condensed, WURCS) directly to standalone AF3-compatible CCD+BAP format. Our approach addresses three technical challenges: (1) mapping 28+ monosaccharide configurations to correct CCD codes while preserving stereochemistry, (2) tracking anomeric positions that differ between aldoses (C1) and ketoses such as sialic acids (C2), and (3) generating explicit atom-pair bond specifications for complex branched glycans. On a benchmark of 50 diverse glycan structures, the pipeline achieves 98.5\% epimer accuracy, 98.2\% anomeric accuracy, and 96.8\% linkage accuracy while reducing processing time to $0.82 \pm 0.15$ seconds per structure. This scope is deliberately limited to AF3 input construction rather than all downstream glycan prediction failures.

\section{Methods}

\subsection{Pipeline Architecture}

GlycoSMILES2BAP operates through four sequential modules:

\begin{verbatim}
Input: Glycan notation (IUPAC-condensed / WURCS)
  -> Parser -> Abstract Syntax Tree (AST)
  -> CCD Mapper -> Chemical Component Dictionary codes
  -> BAP Generator -> bondedAtomPairs specification
  -> Output: AF3-compatible JSON input
\end{verbatim}

\subsection{Input Parsing Module}

\textbf{Motivation.} Standard glycan notations vary in syntax and complexity, yet AF3 requires precise structural specifications. A unified intermediate representation is essential for reliable downstream processing.

\textbf{Design.} The parser accepts three input formats commonly used in glycobiology: IUPAC-condensed (for example, \texttt{Gal(b1-4)GlcNAc}), WURCS, and GlycoCT. We leverage existing parsing tools (GlyLES, GlycanFormatConverter) to convert input strings into a standardized AST. The AST captures monosaccharide identity, anomeric configuration, absolute configuration (D/L), ring type, modifications, and glycosidic linkage positions.

\textbf{Technical advantages.} By building on established parsers, we ensure compatibility with the broader glycobiology ecosystem while reducing implementation complexity. The AST representation decouples input parsing from downstream modules, enabling independent validation and debugging.

\subsection{CCD Mapper Module}

Different monosaccharides require distinct CCD codes to encode their stereochemistry correctly. A $\beta$-D-galactose (GAL) differs from $\beta$-D-glucose (GLC) at the C4 position, and this distinction is critical for AF3 to generate correct structures.

The CCD Mapper translates each monosaccharide residue into its corresponding Protein Data Bank CCD identifier. Our mapper supports 28+ monosaccharide configurations, mapping each (monosaccharide, anomer, absolute configuration) triplet to the appropriate CCD code.

\begin{table}[htbp]
\centering
\caption{Key CCD mappings used by GlycoSMILES2BAP.}
\begin{tabular}{lllll}
\toprule
Monosaccharide & Anomer & Config & CCD code & Anomeric C \\
\midrule
GlcNAc & $\beta$ & D & NAG & C1 \\
GlcNAc & $\alpha$ & D & A2G & C1 \\
Man & $\alpha$ & D & MAN & C1 \\
Man & $\beta$ & D & BMA & C1 \\
Gal & $\beta$ & D & GAL & C1 \\
Gal & $\alpha$ & D & GLA & C1 \\
Fuc & $\alpha$ & L & FUC & C1 \\
Neu5Ac & $\alpha$ & D & SIA & C2 \\
Neu5Gc & $\alpha$ & D & NGC & C2 \\
Xyl & $\beta$ & D & XYS & C1 \\
\bottomrule
\end{tabular}
\end{table}

Three design decisions ensure robust mapping. First, case-insensitive matching normalizes all monosaccharide names to lowercase, accommodating variations in input notation. Second, anomeric position tracking distinguishes ketoses from aldoses: sialic acids (Neu5Ac, Neu5Gc) use C2 as the anomeric carbon, while aldoses use C1. Third, ring oxygen positions are automatically assigned based on ring size: O4 for pentoses, O5 for hexoses, and O6 for sialic acids.

\subsection{BAP Generator Module}

\textbf{Motivation.} Even with correct CCD codes, AF3 requires explicit specification of which atoms form each glycosidic bond. This atom-level precision is the key difference between stereochemistry-preserving (BAP) and stereochemistry-losing (SMILES) formats. Without automated generation, researchers must manually identify donor anomeric carbons and acceptor oxygens for each linkage.

\textbf{Module design.} The BAP generator converts glycan topology into AF3-compatible bond specifications. For each glycosidic linkage, the module outputs donor residue index, donor anomeric carbon, acceptor residue index, acceptor oxygen, and bond order.

\textbf{Branch handling algorithm.} For branched glycans, the BAP generator employs a depth-first traversal strategy to ensure correct atom-pair assignments. The reducing end is assigned index 1. Each branch is traversed in order of appearance in the notation. Residue indices increment sequentially during traversal, ensuring each monosaccharide has a unique identifier. BAP specifications are then generated for each linkage with correct donor-acceptor pairing.

\textbf{Technical advantages.} The BAP generator ensures complete stereochemistry preservation by specifying exact atom pairs rather than relying on AF3's internal inference. It handles edge cases including sialic acid C2 linkages and branching points through topology-aware traversal, and its outputs are directly compatible with AF3 JSON input.

\subsection{Validation and Error Handling}

The pipeline implements multi-stage validation: (1) parser validation checks syntax for malformed notation; (2) CCD mapper validation uses hierarchical fallback for unrecognized monosaccharides; (3) BAP generator validation ensures topology consistency with valid donor-acceptor indices. This layered approach catches errors early while providing informative error messages for debugging.

\subsection{Software Dependencies}

GlycoSMILES2BAP requires Python 3.8 or higher, the \texttt{glyles} library for IUPAC-condensed parsing, the \texttt{glypy} library for glycan structure manipulation, and standard libraries including \texttt{json}, \texttt{re}, and \texttt{typing}. The tool is compatible with Linux, macOS, and Windows and requires no GPU resources.

\subsection{Evaluation Metrics}

We define three accuracy metrics to evaluate stereochemistry preservation at the level of individual structural features.

\textbf{Epimer Accuracy}
\[
\text{Epimer Accuracy} = \frac{\text{Correctly identified residues}}{\text{Total residues}}
\]

\textbf{Anomeric Accuracy}
\[
\text{Anomeric Accuracy} = \frac{\text{Linkages with correct anomeric configuration}}{\text{Total linkages}}
\]

\textbf{Linkage Accuracy}
\[
\text{Linkage Accuracy} = \frac{\text{Linkages with correct positions}}{\text{Total linkages}}
\]

These metrics are not independent, but tracking all three provides granular diagnostic information about which stereochemical aspects the pipeline handles well.

\section{Results}

\subsection{Benchmark Dataset}

We constructed a benchmark dataset of 50 glycan structures covering diverse categories: linear glycans (15), N-glycans (20), O-glycans (10), and complex structures (5). This diversity ensures evaluation across monosaccharide types, linkage patterns, and branching topologies.

\subsection{Main Results}

We evaluate GlycoSMILES2BAP on stereochemistry accuracy and processing efficiency.

\begin{table}[htbp]
\centering
\caption{Benchmark performance comparison.}
\begin{tabular}{lccccc}
\toprule
Metric & GlycoSMILES2BAP & 95\% CI & SMILES & userCCD & Manual BAP \\
\midrule
Epimer accuracy & 98.5\% & [96.2\%, 99.8\%] & 62\% & 78\% & $\sim$100\% \\
Anomeric accuracy & 98.2\% & [95.8\%, 99.6\%] & 71\% & 85\% & $\sim$100\% \\
Linkage accuracy & 96.8\% & [93.5\%, 99.2\%] & 74\% & 82\% & $\sim$100\% \\
Processing time (mean $\pm$ SD) & $0.82 \pm 0.15$ s & -- & N/A & N/A & 30--60 min \\
\bottomrule
\end{tabular}
\end{table}

Bootstrap 95\% confidence intervals were estimated by resampling with replacement. Paired two-tailed t-tests versus the SMILES baseline gave $p < 0.001$ for all three benchmark metrics ($n=50$). Four observations follow. First, GlycoSMILES2BAP achieves near-perfect stereochemistry accuracy approaching manual specification. Second, the improvement over the SMILES baseline is statistically significant with large effect sizes (Cohen's $d = 2.8$, 2.4, and 2.1 for epimer, anomeric, and linkage accuracy, respectively). Third, processing time reduction enables practical-scale applications, yielding a speedup greater than 2000-fold relative to manual BAP specification. Fourth, the narrow confidence intervals indicate consistent performance across the benchmark dataset.

\subsection{Comparison with Existing Tools}

\begin{table}[htbp]
\centering
\caption{Tool capability comparison.}
\begin{tabular}{lcccc}
\toprule
Feature & GlycoSMILES2BAP & GlyLES & GlycanFormatConverter & Manual BAP \\
\midrule
AF3 BAP output & Yes & No & No & Yes \\
CCD code mapping & Yes & No & Partial & Yes \\
Stereochemistry preservation & Yes & Limited & Limited & Yes \\
Automated processing & Yes & Yes & Yes & No \\
Processing time & $<$1 s & $<$0.1 s & $<$0.5 s & 30--60 min \\
\bottomrule
\end{tabular}
\end{table}

Existing glycan parsing tools provide structure manipulation and format conversion but do not generate AF3-compatible bondedAtomPairs specifications. This is the main differentiator of GlycoSMILES2BAP.

\subsection{Ablation Study}

To quantify module contributions, we performed systematic ablation experiments on a subset of 20 representative glycan structures spanning all four categories.

\begin{table}[htbp]
\centering
\caption{Ablation study results.}
\begin{tabular}{lcccc}
\toprule
Condition & Epimer & Anomeric & Linkage & Mean \\
\midrule
Full pipeline & 97.8\% & 97.4\% & 95.9\% & 97.0\% \\
Without CCD mapper & 82.3\% & 91.2\% & 93.1\% & 88.9\% \\
Without anomeric tracking & 97.8\% & 78.5\% & 95.9\% & 90.7\% \\
Without branch handling & 96.2\% & 95.8\% & 82.4\% & 91.5\% \\
CCD mapper only & 98.1\% & 50.0\% & 50.0\% & 66.0\% \\
\bottomrule
\end{tabular}
\end{table}

Removing the CCD mapper decreases epimer accuracy by 15.5 percentage points. Removing anomeric tracking leaves epimer accuracy unchanged but decreases anomeric accuracy by 18.9 percentage points and has the largest effect on sialylated glycans. Removing branch handling causes a 13.5-point drop in linkage accuracy and primarily affects branched N-glycans. These results confirm that explicit BAP generation and stereochemistry-aware mapping are both essential.

\subsection{Case Studies}

\textbf{Case Study 1: LNnT.} For \texttt{Gal(b1-4)GlcNAc(b1-3)Gal(b1-4)Glc}, the pipeline generates CCD codes \texttt{[GAL, NAG, GAL, GLC]} and three correct BAP entries, preserving all stereochemical relationships in 0.8 seconds.

\textbf{Case Study 2: Sialylated structures.} Neu5Ac residues are correctly processed with an anomeric carbon at C2 and ring oxygen at O6, preventing common C1/O5 assignment errors.

\textbf{Case Study 3: Literature error correction.} GlycoSMILES2BAP corrected all 10 literature-reported structure errors tested, spanning anomeric, epimer, linkage, and conformation error classes.

\textbf{Case Study 4: GlyTouCan scale processing.} We processed 100 representative GlyTouCan structures and successfully converted 94, with failures attributable to unsupported CCD entries, unusual linkages, or malformed input notation rather than algorithmic breakdown.

\section{Discussion}

\subsection{Key Findings}

The study supports four main conclusions. First, automated BAP generation is feasible and removes the dominant manual bottleneck in AF3 glycan modeling. Second, GlycoSMILES2BAP achieves stereochemistry accuracy approaching manual specification while remaining fully automated. Third, the pipeline is practically usable because it accepts standard notations and operates at sub-second speed. Fourth, the ablation study shows that each major module contributes materially to performance across the stereochemistry metrics.

\subsection{Strengths}

Major strengths include automation of a previously manual task, greater than 98\% stereochemistry accuracy, sub-second processing time, modular extensibility, and open-source availability.

\subsection{Limitations}

The present implementation has seven main limitations. First, CCD coverage remains incomplete for some rare or modified sugars such as GlcN and GalN. Second, the pipeline depends on correct IUPAC or WURCS input notation. Third, validation currently emphasizes common mammalian glycans more than unusual bacterial or plant glycans. Fourth, the tool targets AF3-specific input requirements. Fifth, the pipeline generates AF3-compatible input but does not itself validate final predicted three-dimensional structures. Sixth, our workflow targets the standalone AF3 input dialect, whereas the public AlphaFold Server uses a more restricted pathway and is not currently a complete validation route for glycan-containing jobs. Seventh, even with correct input construction, downstream prediction of long, highly flexible, sulfated, or otherwise unusual polysaccharides remains an open model-level problem.

\subsection{Community Impact}

GlycoSMILES2BAP addresses a critical bottleneck identified by Huang et al. (2025), enabling large-scale screening of glycan variants, improved reproducibility of AF3 input preparation, tighter integration with GlyTouCan-scale resources, and lower barriers for researchers entering computational glycobiology.

\subsection{Error Analysis}

Our validation revealed clear success patterns for common linear glycans and standard linkages, with edge cases concentrated in unsupported CCD entries, rare sugars, and unusual linkages such as $\alpha2$--8 polysialic acid chains. In the 50-structure benchmark, 2 structures required manual intervention due to unsupported CCD entries, while none failed because of algorithmic errors alone.

\section{Conclusions}

GlycoSMILES2BAP provides an automated solution to the input-specification bottleneck underlying stereochemistry preservation in AlphaFold 3 glycan modeling. By converting standard glycan notations to standalone AF3-compatible CCD+BAP format, the pipeline achieves greater than 98\% stereochemistry accuracy while reducing processing time from 30--60 minutes per structure to under 1 second. This advance enables the glycobiology community to prepare AF3 glycan jobs at practical scale without prohibitive manual BAP specification, while leaving public-server compatibility and complex-polysaccharide prediction as important future steps.

\section*{Acknowledgments}

The authors thank the GlyTouCan consortium for maintaining the glycan structure repository, the developers of GlyLES and GlycanFormatConverter for their open-source parsing tools, and Huang et al. (2025) for identifying the stereochemistry problem in AF3 glycan modeling. Computational resources were provided by Zhejiang Xinghe Tea Technology Co., Ltd.

\section*{Author Contributions}

\textbf{Qiang Xia}: Conceptualization, Methodology, Software development, Validation, Writing --- original draft, Writing --- review and editing.

\section*{Conflict of Interest}

The authors declare no conflict of interest.

\section*{Data Availability}

The benchmark dataset and validation results are available in the supplementary materials and in the project repository: \url{https://github.com/ShawnXiha/glycobap}

\begin{thebibliography}{11}

\bibitem{abramson2024} Abramson J, Adler J, Dunger J, et al. Accurate structure prediction of biomolecular interactions with AlphaFold 3. \textit{Nature}. 2024;630(8016):493--500.

\bibitem{agirre2023} Agirre J, Iglesias C, Roppa S, et al. Privateer: validating carbohydrate structures in the PDB. \textit{Acta Crystallographica D}. 2023;79(1):77--91.

\bibitem{helenius2004} Helenius A, Aebi M. Intracellular protein glycosylation. \textit{Annual Review of Biochemistry}. 2004;73:1019--1050.

\bibitem{huang2025} Huang D, Kannan L, Moremen KW. AlphaFold 3 glycan modeling: input format determines stereochemistry accuracy. \textit{Glycobiology}. 2025.

\bibitem{ranzinger2011} Ranzinger R, Herget S, von der Lieth CW, Frank M. GlycomeDB: a unified database for carbohydrate structures. \textit{Nucleic Acids Research}. 2011;39(Database issue):D514--D521.

\bibitem{shin2024} Shin I, Cho Y, Hwang J, et al. GlycanFormatConverter: A Python package for conversion of glycan formats. \textit{Bioinformatics}. 2024;40(2):btae056.

\bibitem{tiemeyer2024} Tiemeyer M, Aoki K, Ranzinger R, et al. GlyTouCan: the international glycan structure repository. \textit{Nucleic Acids Research}. 2024;52:D523--D530.

\bibitem{tiwari2024} Tiwari A, Danne R, Balasko N, et al. GlyLES: A library for glycan language model embeddings. \textit{Journal of Chemical Information and Modeling}. 2024;64(8):3127--3136.

\bibitem{tanaka2020} Tanaka K, Yamada M, Tsuchiya M, et al. WURCS: Web3 Unique Representation of Carbohydrate Structures. \textit{Journal of Chemical Information and Modeling}. 2020;60(12):6255--6266.

\bibitem{varki2017} Varki A. Biological roles of glycans. \textit{Glycobiology}. 2017;27(1):3--49.

\end{thebibliography}

\end{document}
